\chapter{Fazit}

\section{Zusammenfassung}

Im Projekt wurde ein vollständiger Workflow für Stereo-Bildverarbeitung und
Machine-Learning-basierte Objekterkennung umgesetzt. Der Bildverarbeitungsteil umfasst die
ChArUco-basierte Mono- und Stereo-Kalibrierung, die Speicherung der Kalibrierparameter, die
Aufnahme und Wiederverwendung von Stereo-Bildpaaren, die Rektifizierung, die Berechnung von
Disparitätskarten mit StereoSGBM und den Export farbiger Punktwolken. Ergänzend wurden
Filter zur Entfernung unplausibler Tiefen, Disparitätsartefakte und kleiner getrennter
Komponenten erstellt.

Der Machine-Learning-Teil umfasst einen eigenen YOLO-Datensatz für drei Schreibgeräteklassen
und ein trainiertes YOLO11m-Modell. Der beste Trainingsstand erreicht eine mAP@50 von 0,985
und eine mAP@50:95 von 0,845. Die Inferenz wurde in das Rekonstruktionsnotebook integriert
und erzeugt annotierte linke und rechte Rohbilder.

Wichtig ist die klare Rollenverteilung: Die klassische Bildverarbeitung liefert metrische
3D-Geometrie, das Machine Learning liefert semantische Objektinformation. Beide Teile sind
getrennt nachvollziehbar und bilden gemeinsam eine solide Basis für semantische
3D-Szenenanalyse.

\section{Grenzen}

Die Punktwolkenqualität hängt stark von der Kalibrierung, der Szene, der Beleuchtung und dem
Disparitätsbereich ab. Schwach texturierte oder spiegelnde Oberflächen können weiterhin zu
Fehlmatches führen. Der Stereo-RMS-Wert von ca. $1{,}34\,\mathrm{px}$ ist als
Plausibilitätswert brauchbar, zeigt aber auch, dass eine genauere Kalibrierkampagne die
3D-Qualität verbessern könnte.

Beim YOLO-Modell ist die Validierungsleistung hoch, der Datensatz ist jedoch klein und aus
ähnlichen Aufnahmebedingungen entstanden. Deshalb sollte die gute mAP nicht als Garantie für
beliebige neue Szenen verstanden werden. Besonders andere Hintergründe, Beleuchtungen,
Verdeckungen und Kamerawinkel sollten für eine robuste Bewertung zusätzlich getestet werden.

\section{Ausblick}

Die nächsten sinnvollen Schritte sind:

\begin{itemize}
    \item quantitative Bewertung der 3D-Rekonstruktion mit Referenzobjekten,
    \item zusätzliche Kalibrierbilder mit stärkerer Abdeckung des Bildfeldes,
    \item systematische Parameterstudie für StereoSGBM,
    \item Erweiterung des YOLO-Datensatzes um mehr Hintergründe, Abstände und Verdeckungen,
    \item Auswertung des Testsets zusätzlich zur Validierung,
    \item Fusion von 2D-Bounding-Boxes mit den rekonstruierten 3D-Punkten,
    \item Berechnung objektbezogener 3D-Merkmale wie Schwerpunkt, Ausdehnung und Abstand.
\end{itemize}

\section{Verwendete Quellen}

Als fachliche Grundlage wurden zuerst die bereitgestellten Kursunterlagen verdichtet und auf
das Projekt bezogen:

\begin{itemize}
    \item Stefan Reinmüller: \textit{Industrielle Bildaufnahme \& Bildverarbeitung},
          veröffentlicht am 13.04.2026, \url{https://stefan-reinmueller.github.io/MECH-M-DUAL-2-BEV/}.
    \item Peter Kandolf: \textit{Machine Learning in Industrial Image Processing},
          veröffentlicht am 04.04.2026, \url{https://kandolfp.github.io/MECH-M-DUAL-2-MLB/}.
\end{itemize}

Zusätzlich wurden die Projektdateien, Notebooks, gespeicherten Kalibrierparameter,
YOLO-Konfigurationen, Trainingsmetriken und erzeugten Artefakte im Workspace als direkte
Quellen für die Dokumentation verwendet.
